%! Fundamental Theorem of Algebra and Complex Zeros — Unit 3, Lesson 3.5 — Guided Notes
\documentclass[10pt]{article}
\usepackage{algebra2-article}
\usepackage{algebra2-boxes}
\newcommand{\TallMath}[1]{$\displaystyle #1\rule[-1.4em]{0pt}{3.2em}$}

\newcommand{\lgrid}[4]{%
  \draw[step=1, linegray!40, very thin] (#1,#3) grid (#2,#4);
  \draw[->, semithick, charcoal] (#1,0)--(#2,0) node[below right, font=\scriptsize]{$x$};
  \draw[->, semithick, charcoal] (0,#3)--(0,#4) node[left, font=\scriptsize]{$y$};}

\begin{document}

\pageheader{Unit 3, Lesson 3.5}{Guided Notes}
\namedateperiod

\vspace{0.08in}

\begin{objectivebox}
\textbf{By the end of this lesson, I will be able to\dots}
\begin{itemize}\setlength{\itemsep}{2pt}
  \item state the \emph{Fundamental Theorem of Algebra} and use it to say how many solutions a polynomial equation has \emph{before} solving it;
  \item explain why imaginary zeros of a polynomial with real coefficients come in \emph{conjugate pairs};
  \item determine the \emph{number and type} of solutions (real or imaginary) from a degree and a graph;
  \item solve a polynomial equation of degree $3$ or higher over the \emph{complex} numbers, and write a polynomial from a given list of zeros.
\end{itemize}
\end{objectivebox}

\vspace{0.05in}

\begin{vocabbox}
\small
Fill in each term as we build it together.
\vspace{2pt}
\termblanklong{Fundamental Theorem of Algebra}
\termblanklong{Zeros counted with multiplicity}
\termblanklong{Complex Conjugate Root Theorem}
\termblanklong{Complex conjugate pair}
\termblanklong{Number and type of solutions}
\end{vocabbox}

\begin{hookbox}
Here is $f(x) = x^3 - x^2 + 4x - 4$, the polynomial you finished in the Warm-Up.
(The $y$-axis is compressed: each horizontal gridline is $8$ units.)
\begin{center}
\begin{tikzpicture}[scale=0.5]
  \lgrid{-2}{3}{-4}{4}
  \begin{scope}
    \clip (-2,-4) rectangle (3,4);
    \draw[very thick, forest, domain=-1.95:3.0, samples=150, smooth]
      plot (\x,{((\x)*(\x)*(\x) - (\x)*(\x) + 4*(\x) - 4)/8});
  \end{scope}
  \foreach \x in {-1,1,2}
    \draw[charcoal] (\x,-0.15)--(\x,0.15) node[below=1pt, font=\tiny]{$\x$};
  \fill[charcoal] (1,0) circle (2.5pt);
\end{tikzpicture}
\end{center}
\begin{itemize}\setlength{\itemsep}{1pt}
  \item How many times does this graph cross the $x$-axis? \blank{1.4cm} \; At $x = $ \blank{1.2cm}
  \item How many solutions did you find in the Warm-Up? \blank{1.2cm} \; List them: \blank{3.4cm}
  \item So two of your three solutions are \emph{not} on this graph. What kind of numbers are they?
        \blank{3.0cm}
\end{itemize}
\textbf{Today's question: how many solutions does a polynomial equation \emph{really} have --- and
how many of them can a graph show you?}
\end{hookbox}

\begin{notesbox}{1. The Fundamental Theorem of Algebra}
\begin{tcolorbox}[colback=white, colframe=navy, boxrule=0.7pt, arc=1.5mm,
    left=2mm, right=2mm, top=1.5mm, bottom=1.5mm]
\textbf{Fundamental Theorem of Algebra.} Every polynomial of degree $n \ge 1$ has at least
\blank{1.2cm} \emph{complex} zero.

\vspace{4pt}
\textbf{The corollary we actually use.} A polynomial of degree $n$ has \emph{exactly}
\blank{1.2cm} complex zeros, counted with \textbf{multiplicity}, and factors completely as
\[ f(x) \;=\; a(x - r_1)(x - r_2)\cdots(x - r_n). \]
\end{tcolorbox}

\vspace{4pt}
\textbf{Why the corollary follows.} The theorem promises one zero $r_1$. By the Factor Theorem
(Lesson 3.3), $(x - r_1)$ is then a \blank{2.0cm} of $f$, so dividing leaves a quotient of degree
\blank{1.6cm}. Apply the theorem to \emph{that} polynomial, and repeat. Each pass uses up one degree,
so the process runs exactly \blank{1.2cm} times.

\vspace{5pt}
\textbf{Read that again, because it is new.} Every other tool this year found solutions. This one
\emph{counts} them --- before you find any. Two words matter:
\begin{itemize}[leftmargin=1.5em, itemsep=3pt]
  \item \textbf{complex}: real numbers \emph{are} complex numbers ($5 = 5 + 0i$), so the count
        includes both kinds. It does \emph{not} mean all $n$ are imaginary.
  \item \textbf{counted with multiplicity}: a repeated zero is counted as many times as it
        \blank{3.0cm}.
\end{itemize}
\end{notesbox}

\begin{notesbox}{2. Counting With Multiplicity}
Here is the unit's anchor polynomial $g(x) = x^3 - 3x + 2 = (x-1)^2(x+2)$, whose graph you read back
in Lesson 3.0. (Each horizontal gridline is $2$ units.)
\begin{center}
\begin{tikzpicture}[scale=0.5]
  \lgrid{-3}{3}{-4}{4}
  \begin{scope}
    \clip (-3,-4) rectangle (3,4);
    \draw[very thick, forest, domain=-2.55:2.2, samples=150, smooth]
      plot (\x,{((\x)*(\x)*(\x) - 3*(\x) + 2)/2});
  \end{scope}
  \foreach \x in {-2,-1,1,2}
    \draw[charcoal] (\x,-0.15)--(\x,0.15) node[below=1pt, font=\tiny]{$\x$};
  \fill[charcoal] (-2,0) circle (2.5pt);
  \fill[charcoal] (1,0)  circle (2.5pt);
\end{tikzpicture}
\end{center}
\begin{itemize}[leftmargin=1.5em, itemsep=3pt]
  \item How many \emph{different} numbers are zeros of $g$? \blank{1.2cm} \; Which ones?
        \blank{2.6cm}
  \item The degree is \blank{1.0cm}, so the theorem promises \blank{1.0cm} zeros. That does not match
        --- unless we count $x = 1$ \blank{2.6cm} times, because its factor is
        \blank{2.0cm}.
  \item Zeros of $g$ counted with multiplicity: \blank{4.0cm}
  \item Lesson 3.0 check: at the double zero the graph \blank{2.0cm} the axis; at the single zero it
        \blank{2.0cm} it. The graph shows the multiplicity.
\end{itemize}

\vspace{3pt}
\textbf{The sentence to remember:} \emph{distinct} zeros can be fewer than $n$. Zeros \emph{counted
with multiplicity} are never fewer and never \blank{1.6cm} --- always exactly $n$.
\end{notesbox}

\begin{notesbox}{3. Imaginary Zeros Travel in Pairs}
Look back at the Warm-Up. Item 2(b) gave $1+2i$ and $1-2i$. Item 3 gave $2i$ and $-2i$. Not a
coincidence.

\vspace{4pt}
\begin{tcolorbox}[colback=white, colframe=navy, boxrule=0.7pt, arc=1.5mm,
    left=2mm, right=2mm, top=1.5mm, bottom=1.5mm]
\textbf{Complex Conjugate Root Theorem.} If a polynomial has \emph{\blank{2.0cm}} coefficients and
$a + bi$ (with $b \ne 0$) is a zero, then \blank{2.0cm} is also a zero.
\end{tcolorbox}

\vspace{4pt}
\textbf{Why the pair is necessary.} Multiply the two factors that a conjugate pair contributes:
\[ \big(x - (a+bi)\big)\big(x - (a-bi)\big) \;=\; x^2 - \blank{1.2cm}\,x + \blank{2.0cm} \]
Every $i$ cancels, so the pair contributes a factor with \blank{2.0cm} coefficients. A \emph{lone}
imaginary zero would leave an $i$ stranded in the coefficients --- which is why it can never happen in
a real polynomial.

\vspace{5pt}
\textbf{Two consequences worth more than the theorem itself.}
\begin{itemize}[leftmargin=1.5em, itemsep=4pt]
  \item The number of imaginary zeros is always \blank{2.0cm} ($0$, $2$, $4$, \dots), because they are
        counted in pairs.
  \item Therefore a polynomial of \textbf{odd} degree with real coefficients must have at least
        \blank{1.4cm} real zero. (Pair up the imaginary ones; an odd number cannot be left with
        \blank{2.0cm} real zeros.)
  \item \textbf{Lesson 3.0 agrees.} An odd-degree graph has arms pointing in \blank{2.4cm}
        directions, so it has no choice but to \blank{2.0cm} the $x$-axis at least once.
\end{itemize}
\end{notesbox}

\begin{notesbox}{4. Number and Type --- Without Solving Anything}
An $x$-intercept is a \textbf{real} zero. An imaginary zero is nowhere on the graph. So a graph is an
honest witness, but an \blank{2.6cm} one --- and that gap is exactly what you can now measure:
\[ \text{(number of imaginary zeros)} \;=\; \text{degree} \;-\; \text{(number of \blank{2.4cm} zeros, counted with multiplicity)} \]

\vspace{3pt}
Each row below describes a polynomial with \emph{real} coefficients. Fill in the rest.
\begin{center}
\renewcommand{\arraystretch}{1.5}
\begin{tabularx}{\linewidth}{|l|C{2.2cm}|C{3.0cm}|C{2.6cm}|X|}
\hline
\rowcolor{forestbg}
& \textbf{Degree} & \textbf{Real zeros (with mult.)} & \textbf{Imaginary zeros} & \textbf{How you know} \\ \hline
(a) $f$ from the Hook & $3$ & & & \\ \hline
(b) a quartic whose graph crosses twice & $4$ & $2$ & & \\ \hline
(c) $g(x)=(x-1)^2(x+2)$ & & & & \\ \hline
(d) a quintic whose graph crosses once & $5$ & & & \\ \hline
\end{tabularx}
\end{center}

\vspace{3pt}
\textbf{The question to test yourself with:} could a \emph{cubic} with real coefficients have
\emph{no} $x$-intercepts at all? \blank{1.2cm} \; Why not? \writeline
\end{notesbox}

\begin{notesbox}{5. Solving Over the Complex Numbers}
The workflow is Lesson 3.4's, unchanged: \textbf{list} $\rightarrow$ \textbf{test} $\rightarrow$
\textbf{divide} $\rightarrow$ \textbf{finish}. Only the last step changes --- when the depressed
quadratic has a negative discriminant you no longer write ``no solution.'' You write the
\blank{3.0cm}.

\vspace{5pt}
\textbf{(a) Together:} Solve $x^3 - 5x^2 + 9x - 5 = 0$.

\vspace{2pt}
Candidates: \blank{3.0cm} \quad One that works: $r = $ \blank{1.0cm}
\begin{center}
\renewcommand{\arraystretch}{1.6}
\begin{tabular}{r|C{1.4cm}C{1.4cm}C{1.4cm}C{1.4cm}}
    & $1$ & $-5$ & $9$ & $-5$ \\
    &     &      &     &      \\ \cline{2-5}
    &     &      &     &      \\
\end{tabular}
\end{center}
Depressed polynomial: \blank{3.0cm} \; Discriminant: \blank{1.6cm} \; so the other two solutions are
\blank{2.6cm}.

\vspace{2pt}
Finish it: $x = $ \blank{3.4cm} \qquad All solutions: \blank{4.6cm}

\vspace{2pt}
\textbf{Number and type:} \blank{1.0cm} real and \blank{1.0cm} imaginary --- and
$\blank{0.8cm} + \blank{0.8cm} = 3$, the degree. That check \emph{is} the verification step
(A2.EI.6d).

\vspace{6pt}
\textbf{(b) No division needed.} Solve $x^4 - 5x^2 - 36 = 0$ using \emph{quadratic form}
(Lesson 3.2).

\vspace{2pt}
Factor: $x^4 - 5x^2 - 36 = ($\blank{2.0cm}$)($\blank{2.0cm}$)$

\vspace{2pt}
Set each factor to zero: $x = $ \blank{2.0cm} \quad and \quad $x = $ \blank{2.0cm}

\vspace{2pt}
Four solutions for a degree-\blank{0.8cm} equation: \blank{4.0cm} \; ---
\blank{1.0cm} real, \blank{1.0cm} imaginary.

\vspace{3pt}
\textbf{Note:} over the \emph{real} numbers, $x^2+4$ is as far as you can factor --- it is
irreducible. Over the \emph{complex} numbers it splits into $(x-2i)(x+2i)$, which is what ``factors
completely'' means once imaginary numbers are allowed.
\end{notesbox}

\begin{notesbox}{6. Running It Backward: Building a Polynomial From Its Zeros}
Lesson 3.4 went polynomial $\rightarrow$ zeros. The standard also asks for the reverse: zeros
$\rightarrow$ \blank{2.0cm} $\rightarrow$ polynomial. Each zero $r$ contributes the factor
\blank{1.6cm}.

\vspace{5pt}
\textbf{(a) Real zeros $-1$ (multiplicity 2) and $4$.}

\vspace{2pt}
Factored form: $f(x) = $ \blank{4.0cm} \quad Degree: \blank{1.0cm}

\vspace{2pt}
Standard form (multiply it out): \blank{6.0cm}

\vspace{6pt}
\textbf{(b) Zeros $2$ and $3i$, with \emph{real} coefficients.} \; Before anything else:
which zero must you \emph{also} include, and why? \blank{5.4cm}

\vspace{2pt}
Factored form using the conjugate pair's real quadratic: $f(x) = ($\blank{1.8cm}$)($\blank{2.2cm}$)$

\vspace{2pt}
Standard form: $f(x) = $ \blank{6.0cm}

\vspace{2pt}
Check: the degree is \blank{0.8cm}, and you listed \blank{0.8cm} zeros. \checkmark

\vspace{6pt}
\textbf{(c) The question that tests whether you believe the theorem.} What is the \emph{smallest}
possible degree of a polynomial with real coefficients whose zeros include $5$ and $1 - i$?
\blank{1.2cm} \; Explain: \writeline
\end{notesbox}

\begin{practicebox}
\textbf{A. Solve over $\mathbb{C}$.} $x^3 + 3x^2 + 9x + 27 = 0$ \; (try factoring by grouping first).
\begin{enumerate}[leftmargin=1.7em, itemsep=6pt]
  \item Factored form: \blank{5.0cm}
  \item All solutions: \blank{4.6cm}
  \item Number and type: \blank{1.0cm} real, \blank{1.0cm} imaginary. Does that total match the
        degree? \blank{1.2cm}
\end{enumerate}

\vspace{4pt}
\textbf{B. Count without solving.} A polynomial with real coefficients has degree $6$, and its graph
crosses the $x$-axis at exactly $x = -3$ and $x = 1$ and touches it at $x = 4$.
\begin{enumerate}[leftmargin=1.7em, itemsep=6pt, start=4]
  \item Real zeros counted with multiplicity: \blank{1.4cm} \; (remember what a \emph{touch} means)
  \item Imaginary zeros: \blank{1.2cm} \; How do you know that number had to be even? \blank{4.6cm}
\end{enumerate}

\vspace{4pt}
\textbf{C. Build one.} Write a polynomial of least degree with real coefficients whose zeros are
$-2$ and $1 + 3i$.
\begin{enumerate}[leftmargin=1.7em, itemsep=6pt, start=6]
  \item The full list of zeros: \blank{4.0cm}
  \item Factored form: \blank{5.4cm}
  \item Standard form: \blank{6.0cm}
\end{enumerate}
\end{practicebox}

\end{document}
